\section{Problems}
\begin{pro}
Prove that a number and its square have the same parity.  That is, the square of an even number
is even and the square of an odd number is odd.
\end{pro}

\begin{pro}\label{pro:invConc}
Prove that the inverse of an implication is true if and only if the converse of the implication is true.
\end{pro}

% Eval example now.
%\begin{pro}\label{pro:abeven}
%Prove that if $a$ and $b$ are integers and $ab$ is even, 
%then at least one of $a$ or $b$ is even.
%\end{pro}

\begin{pro}
Let $a$ and $b$ be integers.
Consider the problem of proving that if at least one of $a$ or $b$ is even, then $ab$ is even.
Is this equivalent to the statement from Evaluate~\ref{eval:abeven}?  Explain, using the appropriate
terminology from this chapter.
\end{pro}

\begin{pro}
Let $a$ and $b$ be integers. 
Consider the statement
``{\em If $ab$ is even, then at least one of $a$ or $b$ is even.}''
Rephrase this statement using the word {\em odd} instead of even (but you cannot use the phrase {\em not odd}). 
Using terminology from this chapter, how did you come up with the alternative phrasing?
%Rephrase the statement from Evaluate~\ref{eval:abeven} 
%without using the words {\em even} or {\em not}.
%Using terminology from this chapter, how did you come up with the alternative phrasing?
\end{pro}

\begin{pro}
Prove or disprove that there are 100 consecutive positive integers that are not perfect squares.
(Recall: a number is a perfect square if it can be written as $a^2$ for some integer $a$.)
\end{pro}

\begin{pro}
Consider the equation $n^4+m^4=625$.  
\begin{enumerate}[(a)]
\item
Are there any {\em integers} $n$ and $m$ that satisfy this equation? Prove it.
\item
Are there any {\em positive integers} $n$ and $m$ that satisfy this equation? Prove it.
\end{enumerate}
\end{pro}

\begin{pro}
Consider the equation $a^3+b^3=c^3$ over the integers (that is, $a$, $b$, and $c$ have to all be integers).
\begin{lenum}
\item
Prove that the equations has infinitely many solutions.
\item
If we restrict $a$, $b$, and $c$ to the positive integers, are there infinitely many solutions?  Are there any?
Justify your answer.
(Hint: Do a web search for ``Fermat's Last Theorem.'')
\end{lenum}
\end{pro}


%\begin{pro}
%Let $n$ be an odd integer and $k$ an integer.
%Prove that $kn$ is odd iff $k$ is odd.
%\end{pro}

% Part 2 of the following is tricky. The answer depends on whether they assume
% $x$ is an integer.  I am not sure if I like this ambiguity, but I'll leave it for now.
\begin{pro}
Let $n$ be an integer.
\begin{lenum}
\item
Prove that if $n$ is odd, then $3n+4$ is odd.
\item
Is it possible to prove that $n$ is odd iff $3n+4$ is odd?  If so, prove it.  
If not, explain why not (i.e. give a counter example).
\item
If we don't assume $n$ has to be an integer, 
is it possible to prove that $n$ is odd iff $3n+4$ is odd?  If so, prove it.  
If not, explain why not (i.e. give a counter example).
\end{lenum}
\end{pro}


\begin{pro}
Prove that if $n$ is an integer and $5 n + 4$ is even, then $n$ is
even using a
\begin{lenum}
\item
direct proof
\item
proof by contraposition
\item
proof by contradiction
 \end{lenum}
\end{pro}

% Intentional repeat of the first problem, just rephrased.
\begin{pro}
Prove that $a$ is even if and only if $a^2$ is even.
\end{pro}

\begin{pro}
Prove that $n^2+2n+1$ is even if and only if $n$ is odd.
\end{pro}

\begin{pro}
Let $n$ be an integer.
\begin{lenum}
\item
Prove that if $n$ is odd, then $4n+3$ is odd.
\item
Is it possible to prove that $n$ is odd iff $4n+3$ is odd?  If so, prove it.  
If not, explain why not (i.e. give a counter example).
\end{lenum}
\end{pro}



\begin{pro}
Prove that $ab$ is odd iff $a$ and $b$ are both odd.
\end{pro}

\begin{pro}
Prove or disprove each of the following.
\begin{lenum}
\item
Let $k$ be an odd integer. Then $a$ is even if and only if $k a$ is even.
\item
Let $k$ be an even integer. Then $a$ is even if and only if $k a$ is even.
\item
Let $k$ be an integer. Then $a$ is even if and only if $k a$ is even.
\end{lenum}
\end{pro}

\begin{pro}
Let $n$ be an odd integer. For what values of $k$ do $n$ and $n k$ have the same parity?
Prove your claim.
\end{pro}

\begin{pro}
Let $n$ be an even integer. For what values of $k$ do $n$ and $n k$ have the same parity?
Prove your claim.
\end{pro}

\begin{pro}
Prove or disprove: Every positive integer can be written as the
sum of the squares of two integers.
\end{pro}

\begin{pro}
Prove that the product of two rational numbers is rational.
\end{pro}

\begin{pro}
Prove that the product of a non-zero rational number and an irrational number is
irrational.
\end{pro}

\begin{pro}
Prove or disprove that $c$ is irrational if and only if $c+1$ is irrational. 
\end{pro}

\begin{pro}
Prove or disprove that $c$ is rational if and only if $c^2$ is rational. 
\end{pro}

\begin{pro}
Prove or disprove that $n^2-1$ is composite whenever $n$ is a positive
integer greater than or equal to 1.
\end{pro}

\begin{pro}
Prove or disprove that $n^2-1$ is composite whenever $n$ is a positive
integer greater than or equal to 3. 
\end{pro}

\begin{pro}
Compute $7!$.
\end{pro}

\begin{pro}
Compute $8!/6!$.
\end{pro}

\begin{pro}
List the permutations of the set $\{a,b,c,d\}$.
\end{pro}


\begin{pro}
Prove or disprove that $P=NP$.\footnote{A successful solution to this will earn you an {\em A} in the course.
You are free to use Google or whatever other resources you want for this problem, but you must fully understand
the solution you submit.}
\end{pro}